\documentclass[11pt,a4paper]{article}
\usepackage[utf8]{inputenc}
\usepackage{amsmath,amssymb,amsthm}
\usepackage{listings}
\usepackage{xcolor}
\usepackage[margin=1in]{geometry}
\usepackage{hyperref}

\newtheorem{theorem}{Theorem}
\newtheorem{lemma}[theorem]{Lemma}

\lstset{
  language=C++,
  basicstyle=\ttfamily\small,
  keywordstyle=\color{blue}\bfseries,
  commentstyle=\color{green!60!black},
  stringstyle=\color{red},
  numbers=left,
  numberstyle=\tiny\color{gray},
  breaklines=true,
  frame=single,
  tabsize=4,
  showstringspaces=false
}

\title{IOI 2004 -- Artemis}
\author{Solution}
\date{}

\begin{document}
\maketitle

\section{Problem Statement Summary}
Given $N$ points in the plane, find an axis-aligned rectangle containing
the maximum number of points \emph{strictly} in its interior (no point
on the boundary).

\section{Solution: Sweep with Maximum Subarray}

\subsection{Key Observations}
\begin{enumerate}
  \item Since no point may lie on the boundary, the rectangle's sides
        must be placed in gaps between consecutive distinct coordinate
        values.
  \item Only the $O(N)$ distinct $x$- and $y$-coordinates matter for
        placing sides (coordinate compression).
  \item Fix the left and right boundaries (choosing from gaps between
        consecutive $x$-groups). The interior points form a
        one-dimensional problem: find the $y$-interval maximizing the
        count.
\end{enumerate}

\subsection{Algorithm}
\begin{enumerate}
  \item Group points by their $x$-coordinate and sort the groups.
  \item Compress $y$-coordinates.
  \item For each left boundary $l$ (placed just right of
        $x$-group $l$), sweep the right boundary $r$ rightward:
        \begin{itemize}
          \item When $r$ advances past group $g$, add group $g$'s
                points to a column histogram $\mathrm{cnt}[y]$.
          \item After each update, find the maximum contiguous subarray
                sum of $\mathrm{cnt}[]$ (Kadane's algorithm in $O(M)$
                where $M$ is the number of distinct $y$-values).
        \end{itemize}
  \item The maximum over all $(l, r)$ pairs is the answer.
\end{enumerate}

\subsection{Correctness of Kadane's Step}
Choosing the best $y$-interval corresponds to finding bottom/top
boundaries that maximize the number of interior points. Since boundaries
are placed in gaps between distinct $y$-values, all points in a
contiguous range of compressed $y$-indices are captured. This is
exactly the maximum contiguous subarray sum.

\section{C++ Implementation}

\begin{lstlisting}
#include <bits/stdc++.h>
using namespace std;

int main() {
    ios::sync_with_stdio(false);
    cin.tie(nullptr);

    int N;
    cin >> N;

    vector<int> px(N), py(N);
    for (int i = 0; i < N; i++)
        cin >> px[i] >> py[i];

    // Compress y-coordinates
    vector<int> ys = py;
    sort(ys.begin(), ys.end());
    ys.erase(unique(ys.begin(), ys.end()), ys.end());
    int M = ys.size();

    vector<int> cy(N);
    for (int i = 0; i < N; i++)
        cy[i] = lower_bound(ys.begin(), ys.end(), py[i]) - ys.begin();

    // Sort points by x, group by distinct x-values
    vector<int> order(N);
    iota(order.begin(), order.end(), 0);
    sort(order.begin(), order.end(),
         [&](int a, int b) { return px[a] < px[b]; });

    vector<vector<int>> groups; // groups[g] = list of compressed y-indices
    for (int i = 0; i < N; ) {
        int j = i;
        while (j < N && px[order[j]] == px[order[i]]) j++;
        vector<int> g;
        for (int k = i; k < j; k++)
            g.push_back(cy[order[k]]);
        groups.push_back(g);
        i = j;
    }

    int G = groups.size();
    int ans = 0;

    // Enumerate left boundary (just right of group l)
    for (int l = 0; l < G; l++) {
        vector<int> cnt(M, 0);
        // Right boundary just left of group r; interior = groups l+1..r-1
        for (int r = l + 2; r < G; r++) {
            // Add group r-1 (now strictly inside)
            for (int y : groups[r - 1])
                cnt[y]++;

            // Kadane's algorithm: max contiguous subarray sum of cnt[]
            int best = 0, cur = 0;
            for (int y = 0; y < M; y++) {
                cur += cnt[y];
                best = max(best, cur);
                if (cur < 0) cur = 0;
            }
            ans = max(ans, best);
        }
    }

    cout << ans << "\n";
    return 0;
}
\end{lstlisting}

\section{Complexity Analysis}
\begin{itemize}
  \item \textbf{Time:} $O(G^2 \cdot M)$ where $G \le N$ is the number
        of distinct $x$-groups and $M \le N$ is the number of distinct
        $y$-values. Worst case $O(N^3)$.
  \item \textbf{Space:} $O(N)$.
\end{itemize}

For the IOI constraints ($N \le 20{,}000$), more sophisticated
approaches using balanced BSTs or divide-and-conquer reduce the
complexity to $O(N^2 \log N)$ or $O(N^2)$.

\end{document}
